\section{Background}

This section reviews the key concepts underlying the proposed method -- graphs,
graph-based neural networks, and monocular depth estimation for lifting 2-D
joint coordinates to 2.5-D.

\subsection{Graphs}
Graphs are a versatile structure for representing complex relationships between
entities. A graph consists of nodes (or vertices) and edges (or links) that
connect pairs of nodes, making it particularly suitable for modeling data where
relationships are central, such as the connections between human joints in pose
estimation.

In formal terms, a graph \( G = (V, E) \) includes:
\begin{itemize}
  \item A set of nodes \( V \), representing the entities.
  \item A set of edges \( E \), where each edge \( e \in E \) connects two nodes
  \( (u, v) \), with \( u, v \in V \).
\end{itemize}

Graphs may be directed or undirected, depending on whether the connections
between nodes have a specific direction. This flexibility enables graphs to
effectively capture the various types of relationships that may be beneficial 
for pose estimation.

\subsection{Graph Neural Networks}
Graph Neural Networks (GNNs) \cite{scarselli2008graph} are a class of neural
networks designed to operate on graph-structured data. Unlike traditional neural
networks that work on structured data, GNNs are tailored to handle the irregular
structure of graphs. The fundamental idea behind GNNs is the message passing
framework, which allows nodes to iteratively exchange information with their
neighbors to update their own representations.

In the message passing framework, each node is initially assigned a feature
vector. During each iteration, a node aggregates messages from its neighboring
nodes. This process can be summarized in two main steps:
\begin{itemize}
  \item \textbf{Message Aggregation}: At each iteration \( t \), a node \( v \)
  collects messages from its neighboring nodes \( \mathcal{N}(v) \):
  \begin{equation}
    m_v^{(t+1)} = \mathrm{AGGREGATE}^{(t)} \left( \{ h_u^{(t)} : u \in \mathcal{N}(v) \} \right)
  \end{equation}
  Here, the aggregation function could be a sum, mean, or max function, which
  combines the features of neighboring nodes.
  \item \textbf{Node Update}: After aggregation, the node updates its feature
  vector based on the collected messages:
  \begin{equation}
    h_v^{(t+1)} = \mathrm{UPDATE}^{(t)}(h_v^{(t)}, m_v^{(t+1)})
  \end{equation}
  The update function often involves applying a non-linear transformation, such
  as a neural network layer, to the aggregated messages and the node's current
  state.
\end{itemize}
This iterative process allows each node to refine its representation by
incorporating information from its neighbors, effectively capturing the
dependencies within the graph.

\subsection{Graph Convolutional Networks}
Graph Convolutional Networks (GCNs) \cite{kipf2016semi} are a specialized type
of GNN that extend the idea of convolution from grid-like data structures to
graphs. In a GCN, the convolution operation is adapted to aggregate information
from a node's local neighborhood, enabling the network to learn features that
are informed by the local context within the graph.

A typical GCN layer is defined as:
  \begin{equation}
    H^{(t+1)} = \sigma \left( \tilde{D}^{-1/2} \tilde{A} \tilde{D}^{-1/2} H^{(t)} W^{(t)} \right)
  \end{equation}
where:
\begin{itemize}
  \item \( \tilde{A} = A + I \) is the adjacency matrix of the graph, with added
  self-loops to include the node's own features in the aggregation.
  \item \( \tilde{D} \) is the degree matrix corresponding to \( \tilde{A} \),
  used for normalization.
  \item \( H^{(t)} \) represents the feature matrix at layer \( t \), where each
  row corresponds to a node's features.
  \item \( W^{(t)} \) is the trainable weight matrix for the layer.
  \item \( \sigma \) denotes an activation function, such as ReLU, applied
  element-wise to introduce non-linearity.
\end{itemize}
This layer allows the network to capture the structural information of a node's
neighborhood, leading to more contextually informed node representations.

\subsection{Graph Attention Networks}
Graph Attention Networks (GATs) \cite{velivckovic2017graph} introduce an
attention mechanism (this mechanism is notably used in transformers
\cite{vaswani2017attention} which are the backbone of large language models
\cite{zhao2023survey}) into the GNN framework, allowing the network to assign
different levels of importance to different neighbors during the aggregation
process. This is particularly useful in scenarios where certain neighbors
contribute more significantly to a node's representation than others.

The key steps in a GAT layer are:
\begin{itemize}
  \item \textbf{Self-Attention Mechanism}: For each pair of connected nodes \(
  (u, v) \), an attention coefficient is computed to measure the importance of
  node \( u \)’s features to node \( v \):
  \begin{equation}
    e_{uv} = \mathrm{LeakyReLU} \left( \mathbf{a}^T \left[ W h_u \, || \, W h_v \right] \right)
  \end{equation}
  Here, \( \mathbf{a} \) is a learnable weight vector, \( W \) is a shared
  weight matrix, and \( || \) denotes concatenation.
  \item \textbf{Normalization}: The attention coefficients are normalized across
  all neighbors of a node using the softmax function:
  \begin{equation}
    \alpha_{uv} = \frac{\exp(e_{uv})}{\sum_{k \in \mathcal{N}(v)} \exp(e_{vk})}
  \end{equation}
  This normalization ensures that the sum of the attention weights for each node
  is 1.
  \item \textbf{Aggregation}: Each node’s new feature is computed as a weighted
  sum of its neighbors’ features, with the attention coefficients serving as
  weights:
  \begin{equation}
    h_v' = \sigma \left( \sum_{u \in \mathcal{N}(v)} \alpha_{uv} W h_u \right)
  \end{equation}
  The result is a more nuanced node representation that accounts for the varying
  significance of different neighbors, improving the network's ability to learn
  complex patterns in the data.
\end{itemize}

\subsection{Monocular Depth Estimation} 

Monocular depth estimation infers a dense depth map $D \in \mathbb{R}^{H \times
W}$ from a single RGB image $I \in \mathbb{R}^{H \times W \times 3}$. Because
absolute scale is ambiguous in a single view, networks learn scene priors that
connect image appearance to plausible metric or relative depth. Most
contemporary methods follow an encoder–decoder pattern first popularised by
multi-scale CNNs \cite{eigen2014depth,laina2016deeper}: an encoder extracts
hierarchical features, and a decoder progressively upsamples them—often with
skip connections or multi-scale refinement—to yield per-pixel depth
predictions.

Training regimes fall into two broad categories. \emph{Supervised} approaches
optimise losses such as $\ell_1$, $\ell_2$, or scale-invariant errors against
ground-truth depth captured with active sensors or multiview reconstruction
\cite{laina2016deeper}. \emph{Self-supervised} approaches eliminate depth
labels altogether and instead enforce photometric consistency across stereo
pairs or successive video frames, using differentiable view synthesis and
occlusion masks \cite{godard2017unsupervised,zhou2017unsupervised}.

Recent sota models illustrate the maturation of the field. \textsc{MiDaS}
combines a vision-transformer encoder with a multi-resolution decoder and is
trained on a heterogeneous mixture of data sets to generalise across domains
\cite{ranftl2021robust,ranftl2022dense}. \textsc{ZoeDepth} extends this idea by
mixing real and synthetically rendered images \cite{khattak2023zoedepth}. These
networks now deliver accurate, scale-consistent depth maps with a single
forward pass, making depth a reliable signal for vision tasks.